\begin{center}
\begin{tikzpicture}[x=.88cm,y=.88cm,font=\small]
  \fill[paper] (0,0) rectangle (12.3,6.65);

  \draw[source!70,line width=.8pt] (.5,2.02) .. controls (3.35,1.48) and
    (7.05,2.66) .. (11.85,1.76);
  \node[text=source,below,font=\scriptsize] at (8.2,1.84) {长江方向（示意）};
  \draw[source!55,dashed] (.72,4.46) .. controls (3.7,4.88) and
    (7.4,4.0) .. (11.68,4.45);
  \node[text=source,above,font=\scriptsize] at (8.8,4.27) {淮河方向（示意）};

  \fill[dynasty!44] (6.82,4.58) ellipse (1.42 and .9);
  \node[align=center,text=white] at (6.82,4.58) {洛阳\\西晋中枢};

  \fill[timeline!28] (2.0,2.75) ellipse (1.3 and .9);
  \node[align=center] at (2.0,2.75) {益州、巴蜀\\王濬水军};

  \fill[source!25] (4.62,2.35) ellipse (1.2 and .76);
  \node[align=center] at (4.62,2.35) {荆州、江陵\\杜预军};

  \fill[dynasty!29] (8.86,2.35) ellipse (1.27 and .82);
  \node[align=center] at (8.86,2.35) {江夏、淮南\\诸路晋军};

  \fill[timeline!35] (10.74,.93) ellipse (1.14 and .68);
  \node[align=center,text=white] at (10.74,.93) {建业\\孙吴朝廷};

  \draw[->,very thick,color=dynasty]
    (2.98,2.43) .. controls (4.12,2.0) and (5.48,1.9) .. (6.98,1.88);
  \node[text=dynasty,above,font=\scriptsize,align=center] at (4.9,2.22)
    {益州舰队沿江东下};

  \draw[->,very thick,color=dynasty]
    (5.6,2.2) .. controls (6.72,1.58) and (8.6,1.2) .. (9.67,1.04);
  \node[text=dynasty,below,font=\scriptsize,align=center] at (7.52,1.08)
    {荆州军与水军\\向建业会合};

  \draw[->,thick,dashed,color=timeline]
    (8.05,2.5) .. controls (8.9,2.85) and (9.5,3.33) .. (9.85,3.72);
  \node[text=timeline,right,font=\scriptsize,align=left] at (8.92,3.34)
    {淮南、江夏\\牵制与并进};

  \draw[->,thick,color=source]
    (7.45,4.1) .. controls (7.78,3.6) and (8.0,3.28) .. (8.15,3.03);
  \node[text=source,left,font=\scriptsize,align=right] at (7.34,3.5)
    {北方中枢调度\\但非单一路线};

  \node[draw=source!75,fill=paper,rounded corners=2pt,align=center,
    inner sep=3pt,font=\scriptsize] at (3.44,5.38)
    {280年孙吴降服前后的伐吴把益州水军、荆州军与江淮方向的兵力接入同一攻势；\\
    各军的实际行程、兵额和协同细节仍有异文};
\end{tikzpicture}

\smallskip
\small\textit{图\quad 西晋伐吴的多路水陆走廊，280年：据《晋书·武帝纪》、〈王濬传〉〈杜预传〉〈王浑传〉，《三国志·孙皓传》及《资治通鉴》卷八十至八十一，概括益州、荆州、江夏、淮南与建业之间的部署。城市与河流位置约略；箭线只表示各路行动的概括方向，并不复原舰队、部队、渡口或战斗的精确路线。}
\end{center}

\phantomsection\label{fig:western-jin-unification}
